% Generated draft: requires mathematical and source review.
\documentclass[11pt]{article}
\usepackage[margin=25mm]{geometry}
\usepackage{fontspec}
\setmainfont{Latin Modern Roman}
\usepackage{amsmath,amssymb,mathtools,unicode-math}
\setmathfont{Latin Modern Math}
\usepackage{enumitem,needspace,xurl,hyperref}
\setlength{\parindent}{0pt}
\setlength{\parskip}{5pt}
\setlength{\emergencystretch}{3em}
\newcommand{\N}{\mathbb{N}}
\newcommand{\Z}{\mathbb{Z}}
\newcommand{\Q}{\mathbb{Q}}
\newcommand{\R}{\mathbb{R}}
\newcommand{\C}{\mathbb{C}}
\newcommand{\F}{\mathbb{F}}
\newcommand{\A}{\mathbb{A}}
\newcommand{\PP}{\mathbb{P}}
\newcommand{\E}{\mathbb{E}}
\newcommand{\eps}{\varepsilon}
\newcommand{\dd}{\,\mathrm{d}}
\begin{document}
\section*{UnsolvedMath ID: 700682}
\textbf{Source problem code:} catalogue entry
\subsection{Definitions and mathematical statement}
Let \(p\) be a prime. A group is locally nilpotent if each of its finitely generated subgroups is nilpotent. For a group \(G\), put \(\gamma_1(G)=G\) and \(\gamma_{i+1}(G)=[\gamma_i(G),G]\). For a positive integer \(c\), let \(\mathfrak N_c\) denote the variety defined by \(\gamma_{c+1}(G)=1\), and let \(\mathfrak B_p\) denote the variety defined by \(g^p=1\) for every \(g\in G\); thus the exponent divides \(p\).

Write \(\mathcal V_p=LN\mathfrak M_p\) for the source's class of locally nilpotent groups admitting a splitting automorphism of prime order \(p\), regarded as groups with operators. The source attributes to E. I. Khukhro the assertion that \(\mathcal V_p\) is a variety of groups with operators. The definition of a splitting automorphism is referred to Problem 10.59, whose text is not supplied. Consequently, the defining condition for \(\mathcal V_p\) remains unresolved in this draft.

In a fixed common signature of groups with operators, interpret \(\mathfrak N_c\cap\mathcal V_p\) and \(\mathfrak B_p\cap\mathcal V_p\) by imposing the indicated conditions on the underlying groups. For two varieties \(\mathcal U\) and \(\mathcal W\) in that signature, their join \(\mathcal U\vee\mathcal W\) is the smallest variety containing both.

The proposed question is whether, for every prime \(p\), there exists a positive integer \(c(p)\), depending only on \(p\), such that
\[
\mathcal V_p
=
\bigl(\mathfrak N_{c(p)}\cap\mathcal V_p\bigr)
\vee
\bigl(\mathfrak B_p\cap\mathcal V_p\bigr).
\]
Here the join must be taken in the operator signature intended by the source. That signature and the splitting condition require confirmation from Problem 10.59. There is also an unresolved order convention: requiring the distinguished automorphism to have order exactly \(p\) is not preserved under homomorphic images, since the trivial image has only the identity automorphism. The source's variety assertion therefore requires clarification of how its phrase ``of prime order \(p\)'' is implemented; this draft does not silently replace exact order by order dividing \(p\).
\subsection{Short English statement}
For each prime \(p\), is the stated variety of locally nilpotent groups with a splitting automorphism generated by its members of nilpotency class bounded solely in terms of \(p\), together with its members of exponent dividing \(p\)? The missing definition from Problem 10.59 and the operator order convention must be checked before this becomes a fully specified statement.
\subsection{Sources}
\begin{itemize}
\item Supplied local UnsolvedMath registry record: permanent ID \(700682\), exact source code \texttt{catalogue entry}, title ``Kourovka Notebook Problem 10.67.'' The registry labels the problem open; that label is supplied metadata, not an independently verified assessment.
\item \emph{The Kourovka Notebook (Unsolved Problems in Group Theory)}, Problem 10.67, proposed by E. I. Khukhro in the 10th issue (1986). The supplied background identifies the 21st issue edition (2026), arXiv:1401.0300v48. The collection metadata names E. I. Khukhro and V. D. Mazurov as editors.
\item E. I. Khukhro, \emph{Math. USSR Sbornik} \textbf{58} (1987), 119--126. The supplied statement cites this reference for the assertion that \(LN\mathfrak M_p\) is a variety of groups with operators; no article title is supplied.
\item Kourovka Notebook Problem 10.59 is the supplied cross-reference for the definition of a splitting automorphism. Its text is absent from the supplied record.
\end{itemize}
\end{document}
